% ===== PRÉAMBULE CANONIQUE — à coller VERBATIM en tête de CHAQUE .tex =====
\documentclass[11pt,a4paper]{article}
\usepackage[utf8]{inputenc}\usepackage[T1]{fontenc}\usepackage{lmodern}
\usepackage[french]{babel}
\usepackage{amsmath,amssymb,mathtools}\usepackage{icomma}
\usepackage[version=4]{mhchem}          % \ce{...} : équations & formules
\usepackage{siunitx}\sisetup{locale=FR} % \qty{}{}, \si{}, \ang{}
\usepackage{chemfig}                    % \chemfig{...} : structures développées
\usepackage{xcolor}\usepackage{colortbl}
\usepackage{tikz}\usepackage{pgfplots}\pgfplotsset{compat=1.18}
\usepackage[most]{tcolorbox}\usepackage{enumitem}\usepackage{array,booktabs}
\usepackage[a4paper,margin=2.2cm]{geometry}
\usetikzlibrary{arrows.meta,calc,angles,quotes,positioning,patterns,backgrounds,shapes.geometric}
\definecolor{primary}{RGB}{222,49,99}\definecolor{primarydark}{RGB}{150,22,55}
\definecolor{defcol}{RGB}{0,114,178}\definecolor{methcol}{RGB}{52,92,148}
\definecolor{excol}{RGB}{0,135,90}\definecolor{attncol}{RGB}{200,40,55}
\tcbset{boxbase/.style={enhanced,breakable,boxrule=0.9pt,arc=2.5pt,
  left=3mm,right=3mm,top=2.3mm,bottom=2.3mm,fonttitle=\bfseries,coltitle=white}}
% --- boîtes TUTORIEL (noms EXACTS, ne pas renommer) ---
\newtcolorbox{cle}[1][]{boxbase,colback=defcol!6,colframe=defcol,
  title={\textsc{À retenir}\ifx\relax#1\relax\relax\else\textmd{ : \itshape #1}\fi}}
\newtcolorbox{methode}[1][]{boxbase,colback=methcol!7,colframe=methcol,sharp corners,
  title={\textsc{Méthode}\ifx\relax#1\relax\relax\else\textmd{ --- \itshape #1}\fi}}
\newtcolorbox{exemplebox}[1][]{boxbase,colback=excol!5,colframe=excol,
  title={\textsc{Exemple}\ifx\relax#1\relax\relax\else\textmd{ : \itshape #1}\fi}}
\newtcolorbox{piege}[1][]{boxbase,colback=attncol!6,colframe=attncol,
  title={\textsc{Piège}\ifx\relax#1\relax\relax\else\textmd{ : \itshape #1}\fi}}
\newtcolorbox{remarque}[1][]{boxbase,colback=black!4,colframe=black!45,
  title={\textsc{Remarque}\ifx\relax#1\relax\relax\else\textmd{ : \itshape #1}\fi}}
% --- boîtes EXERCICE / CORRIGÉ (noms EXACTS, auto-comptées) ---
\newtcolorbox[auto counter]{exo}[1][]{boxbase,colback=primary!4,colframe=primary,
  title={\textsc{Exercice}~\thetcbcounter\ifx\relax#1\relax\relax\else\textmd{ --- \itshape #1}\fi}}
\newtcolorbox[auto counter]{corrige}[1][]{boxbase,colback=excol!4,colframe=excol,
  title={\textsc{Corrigé}~\thetcbcounter\ifx\relax#1\relax\relax\else\textmd{ --- \itshape #1}\fi}}
\newcommand{\R}{\mathbb{R}}\newcommand{\N}{\mathbb{N}}\newcommand{\Z}{\mathbb{Z}}\newcommand{\ds}{\displaystyle}
% (un \usepackage{fancyhdr}+en-tête est possible ; il est ignoré côté web)
% =========================================================================
\begin{document}
\sisetup{per-mode=symbol}
\begin{center}\color{primarydark}\large\bfseries Thème 5 — Piles et électrolyse \quad$\bullet$\quad Niveau Difficile\end{center}

\begin{exo}[vrai ou faux : une pile fer--argent]
Pile construite avec les couples \ce{Fe^{2+}}/\ce{Fe} ($E_{\mathrm{r}}^{\circ}=-0{,}44$~V) et
\ce{Ag+}/\ce{Ag} ($E_{\mathrm{r}}^{\circ}=+0{,}80$~V). Vrai ou faux ? Combien de propositions sont
vraies ?
\begin{enumerate}[label=\alph*)]
  \item Le couple \ce{Fe^{2+}}/\ce{Fe} occupe la cathode de la pile.
  \item La f.é.m.\ standard de cette pile vaut $1{,}24$~V.
  \item L'argent métallique est oxydé en ions \ce{Ag+}.
  \item Les électrons circulent de l'électrode de fer vers l'électrode d'argent.
\end{enumerate}
\end{exo}

\begin{exo}[la bonne notation avec des électrodes inertes]
Une pile associe une électrode de platine \ce{Pt} plongée dans \ce{Fe^{3+}}/\ce{Fe^{2+}}
($E_{\mathrm{r}}^{\circ}=+0{,}77$~V) et une électrode de \ce{Pt} plongée dans \ce{Sn^{4+}}/\ce{Sn^{2+}}
($E_{\mathrm{r}}^{\circ}=+0{,}15$~V).
\begin{enumerate}[label=\alph*)]
  \item Quelle demi-pile est la cathode ?
  \item Écris la notation conventionnelle correcte (anode à gauche, électrodes inertes aux extrémités).
\end{enumerate}
\end{exo}

\begin{exo}[électrolyse : de la charge à la masse]
On électrolyse une solution de \ce{CuSO4} avec $I=\qty{2.0}{A}$ pendant $t=\qty{965}{s}$. À la cathode :
\ce{Cu^{2+} + 2 e^- -> Cu} ($M(\ce{Cu})=\qty{63.5}{g\per mol}$, $F=\qty{96485}{C\per mol}$).
\begin{enumerate}[label=\alph*)]
  \item Calcule la charge $Q$ et le nombre de moles d'électrons.
  \item En déduis le nombre de moles de cuivre, puis la masse déposée.
\end{enumerate}
\end{exo}

\begin{exo}[quelle intensité pour argenter ?]
Pour argenter un objet, on veut déposer $m=\qty{3.24}{g}$ d'argent en $t=\qty{1930}{s}$. Données :
\ce{Ag+ + e^- -> Ag} ($n=1$), $M(\ce{Ag})=\qty{108}{g\per mol}$, $F=\qty{96485}{C\per mol}$. Quelle
intensité $I$ faut-il imposer ?
\end{exo}

\begin{exo}[électrolyse de l'eau et volume de gaz]
Lors de l'électrolyse de l'eau, il se forme du dihydrogène à la cathode :
\ce{2 H+ + 2 e^- -> H2}. On impose $I=\qty{9.65}{A}$ pendant $t=\qty{1000}{s}$
($F=\qty{96485}{C\per mol}$ ; volume molaire aux CNTP $=\qty{22.4}{L\per mol}$).
\begin{enumerate}[label=\alph*)]
  \item Calcule la charge $Q$ et le nombre de moles d'électrons.
  \item Combien de moles de \ce{H2} se forment-elles ($2$ électrons par \ce{H2}) ?
  \item Quel volume de \ce{H2} obtient-on aux CNTP ?
\end{enumerate}
\end{exo}

\end{document}
